\documentclass[11pt]{article}
\usepackage[a4paper,margin=1in]{geometry}
\usepackage[T1]{fontenc}
\usepackage{helvet}
\renewcommand{\familydefault}{\sfdefault}
\usepackage{xcolor}
\usepackage{booktabs}
\usepackage{array}
\usepackage{colortbl}
\usepackage{float}
\usepackage[font=footnotesize,labelformat=empty]{caption}
\usepackage{titlesec}
\usepackage{fancyhdr}
\usepackage{enumitem}
\usepackage{amsmath}
\usepackage{lastpage}
\usepackage[hidelinks]{hyperref}

\definecolor{bcnavy}{RGB}{11,31,58}
\definecolor{bcblue}{RGB}{32,74,135}
\definecolor{bcgold}{RGB}{176,141,87}
\definecolor{bcgrey}{RGB}{90,98,112}
\definecolor{bcrule}{RGB}{210,214,220}
\definecolor{bckeep}{RGB}{232,237,244}
\definecolor{bcact}{RGB}{247,241,228}
\definecolor{bcbad}{RGB}{247,230,230}

\titleformat{\section}{\color{bcnavy}\Large\bfseries}{\thesection}{0.6em}{}
\titleformat{\subsection}{\color{bcblue}\large\bfseries}{\thesubsection}{0.6em}{}
\titlespacing*{\section}{0pt}{1.4em}{0.5em}

\pagestyle{fancy}\fancyhf{}
\renewcommand{\headrulewidth}{0pt}
\renewcommand{\footrulewidth}{0.4pt}
\renewcommand{\footrule}{\hbox to\headwidth{\color{bcrule}\leaders\hrule height \footrulewidth\hfill}}
\fancyfoot[L]{\footnotesize\color{bcgrey}Bayesian Capital}
\fancyfoot[C]{\footnotesize\color{bcgrey}Confidential --- Internal Research}
\fancyfoot[R]{\footnotesize\color{bcgrey}Page \thepage\ of \pageref{LastPage}}

\newcommand{\kpi}[2]{\textbf{#1}\,{\footnotesize\color{bcgrey}#2}}
\setlength{\parindent}{0pt}\setlength{\parskip}{0.55em}

\begin{document}

\begin{titlepage}
\vspace*{1.2cm}
{\color{bcgold}\rule{\textwidth}{2pt}}\\[0.4cm]
{\color{bcnavy}\Huge\bfseries Diversifier Candidates}\\[0.35cm]
{\color{bcblue}\LARGE Suitability Assessment}\\[0.5cm]
{\color{bcgold}\rule{\textwidth}{1pt}}\\[0.6cm]
{\large\color{bcgrey} Hybrid Bayesian / Ornstein--Uhlenbeck Volatility Harvester}\\[0.2cm]
{\large\color{bcgrey} OXY \enskip DVN \enskip FSLR \enskip COIN \enskip MSTR \enskip HOOD
 \enskip RBLX \enskip MRNA \enskip ARM \enskip SHOP}\\[1.2cm]
{\color{bcnavy}\large\bfseries Bayesian Capital}\\[0.15cm]
{\color{bcgrey} Quantitative Research}\\[0.15cm]
{\color{bcgrey} 30 July 2026}\\[2cm]

\fbox{\parbox{0.92\textwidth}{\color{bcnavy}\textbf{One-line conclusion.}\\ \color{black}
Assessed as decorrelating additions to the AI-heavy book. A five-fold walk-forward puts
\textbf{DVN and OXY} first (energy; positive in \emph{every} unseen fold), then \textbf{MRNA}
(biotech) and a modest \textbf{FSLR} (solar) --- all decorrelated from the book. \textbf{SHOP}
and \textbf{ARM} add return but little decorrelation; \textbf{COIN} is high-octane, size down.
\textbf{Reject MSTR and HOOD} --- their large in-sample returns collapse to \emph{losses} on
unseen data.}}
\end{titlepage}

% ================= metrics + OOS primer =================
\section{How to read these numbers, and what ``out-of-sample'' means}

Return alone is misleading, so each name is judged on several measures. \emph{Annual return}
is the annualised growth; \emph{Sharpe} is return per unit of volatility (higher $=$ smoother);
\emph{max drawdown} is the worst peak-to-trough equity fall; \emph{Bayes--OU correlation} is the
daily-return correlation between the two sleeves (we want it \emph{low}); \emph{correlation to
the book} is the daily-return correlation to the existing ten names (again, \emph{low} is the
point of the exercise).

\subsection{Out-of-sample testing --- the decisive check}

These candidates arrived with \emph{no} fitted parameters, so we had to fit them. That creates a
trap: an optimiser handed a stretch of history will always find parameters that look brilliant
\emph{on that history}, because it fits the noise as well as the signal. Those in-sample numbers
flatter every strategy --- even a worthless one --- so they cannot tell you whether an edge is
real. Only performance on data the parameters \emph{never saw} can.

\textbf{The mechanics (a train/test holdout).}
\begin{enumerate}[leftmargin=1.6em,itemsep=1pt]
  \item \textbf{Split the timeline.} Fit on the first 60\% of the history; hold out the last 40\%.
  \item \textbf{Fit on train only.} The optimiser sees nothing past the split point.
  \item \textbf{Freeze} the parameters --- no further tuning.
  \item \textbf{Run them forward} on the untouched test window and measure the return there.
        That figure is genuinely out-of-sample: it mimics trading frozen rules into an unknown
        future.
  \item \textbf{Judge on the test window.} If it holds up, the edge is structural. If it
        collapses, the in-sample number was a mirage.
\end{enumerate}

\textbf{Why it is the whole ballgame.} Across this entire engagement, re-optimised parameters
looked superb in-sample and then \emph{lost} out of sample; the frozen parameters always won.
The holdout is the only honest arbiter, and --- as \S5 shows --- it is exactly what separates the
genuine candidates here from the two that merely \emph{look} spectacular.

% ================= scope =================
\section{Scope, data and caveats}

Ten names were supplied as a raw OHLC feed (Apr~2024--Jul~2026, $\sim$584 trading days). Each
was fitted with a robust differential-evolution search (perturbation-averaged objective plus a
minimum-trade floor) and run through the validated engine at \$6M per name; percentage metrics
are scale-invariant.

\textbf{Three requested columns could not be produced from this file.} The feed carries
\emph{OHLC only} --- no volume, so \textbf{average \$ turnover} (volume $\times$ price) cannot be
computed; and \textbf{market cap} and \textbf{P/E} are fundamental data, not derivable from
prices. Supply a volume column and the market-cap / P-E figures (or shares-outstanding and EPS)
and those three fields drop straight in.

\textbf{One methodological note.} Because these parameters are freshly fitted (unlike the earlier
CCL/MU/CVNA/LLY candidates, which arrived pre-fitted), the full-sample metrics are an
\emph{in-sample ceiling}. The out-of-sample holdout (\S5) is the figure to trust.

% ================= volatility + diversification =================
\section{Volatility and diversification}

All ten clear the high-volatility bar. The decisive screen is the right-hand column --- how
correlated each name is to the existing book, since the entire purpose is to dilute the AI/semis
factor that currently dominates it.

\begin{table}[H]\centering\small
\renewcommand{\arraystretch}{1.15}
\begin{tabular}{l l r r c c}
\toprule
\textbf{Name} & \textbf{Sector} & \textbf{Ann vol} & \textbf{Buy-hold DD} & \textbf{corr $\to$ book} & \textbf{corr $\to$ AI-core}\\
\midrule
\rowcolor{bckeep} OXY  & Energy (E\&P)          & 33\% & 48\% & \textbf{0.14} & 0.15\\
\rowcolor{bckeep} DVN  & Energy (E\&P)          & 36\% & 51\% & \textbf{0.21} & 0.21\\
\rowcolor{bckeep} MRNA & Health care (biotech)  & 67\% & 87\% & \textbf{0.27} & 0.22\\
\rowcolor{bckeep} FSLR & Clean energy (solar)   & 59\% & 60\% & 0.32 & 0.33\\
RBLX & Comms / gaming         & 55\% & 71\% & 0.41 & 0.35\\
MSTR & Crypto (BTC treasury)  & 87\% & 83\% & 0.50 & 0.41\\
SHOP & Consumer internet      & 58\% & 47\% & 0.52 & 0.45\\
COIN & Financials (crypto ex.)& 77\% & 66\% & 0.57 & 0.48\\
ARM  & Semiconductors         & 71\% & 54\% & 0.61 & \textbf{0.63}\\
HOOD & Financials (brokerage) & 71\% & 57\% & 0.65 & 0.55\\
\bottomrule
\end{tabular}
\caption*{\footnotesize Blue rows are the genuine decorrelators (energy, biotech, solar). ARM is
the opposite of a diversifier --- at 0.63 to the AI-core it is simply another semiconductor bet.
HOOD/COIN move with risk-on/crypto beta. ``Ann vol'' is annualised realised volatility (a
high-beta proxy; no index was supplied for a literal beta).}
\end{table}

% ================= model fit + OOS =================
\section{Model fit and the out-of-sample holdout}

\begin{table}[H]\centering\small
\renewcommand{\arraystretch}{1.2}
\begin{tabular}{l r r r r c c c c}
\toprule
& \multicolumn{5}{c}{\textbf{Full-sample fit (in-sample ceiling)}} & \multicolumn{2}{c}{\textbf{60/40 holdout}} & \\
\cmidrule(lr){2-6}\cmidrule(lr){7-8}
\textbf{Name} & ann & Sh & maxDD & /yr & B--OU & train & \textbf{test (OOS)} & \textbf{verdict}\\
\midrule
\rowcolor{bckeep} MRNA & 128\% & 2.20 & 28\% & 84  & 0.17 & 62\%  & \textbf{+103\%} & holds\\
\rowcolor{bckeep} OXY  & 48\%  & 2.11 & 10\% & 59  & 0.20 & 27\%  & \textbf{+39\%}  & holds\\
\rowcolor{bckeep} DVN  & 63\%  & 3.48 & 10\% & 125 & 0.20 & 31\%  & \textbf{+13\%}  & holds (modest)\\
\rowcolor{bckeep} FSLR & 94\%  & 1.80 & 24\% & 78  & 0.33 & 117\% & \textbf{+29\%}  & holds\\
SHOP & 107\% & 1.85 & 34\% & 55  & 0.37 & 142\% & +89\% & holds\\
ARM  & 161\% & 2.67 & 22\% & 106 & 0.26 & 173\% & +80\% & holds\\
COIN & 255\% & 3.06 & 41\% & 52  & $-$0.07 & 296\% & +30\% & holds (big fade)\\
RBLX & 140\% & 2.18 & 37\% & 63  & 0.85 & 192\% & +31\% & holds (weak hedge)\\
\rowcolor{bcbad} HOOD & 194\% & 2.57 & 27\% & 74  & 0.35 & 367\% & \textbf{$-$6\%}  & \textbf{FAILS}\\
\rowcolor{bcbad} MSTR & 97\%  & 1.28 & 43\% & 84  & 0.09 & 360\% & \textbf{$-$61\%} & \textbf{FAILS}\\
\bottomrule
\end{tabular}
\caption*{\footnotesize Sorted by suitability. Full-sample columns are the in-sample ceiling;
\textbf{the test column is the honest number}. Blue rows are the recommended diversifiers; red
rows fail out of sample. ``B--OU'' is the internal Bayes--OU hedge correlation (lower is better;
COIN's $-0.07$ is the best we have seen). ``/yr'' is trades per year.}
\end{table}

% ================= overfitting lesson =================
\section{The overfitting lesson --- why the holdout earned its keep}

The two names with the \emph{highest} in-sample fits are the two that fail:

\begin{table}[H]\centering\small
\renewcommand{\arraystretch}{1.15}
\begin{tabular}{l r r r}
\toprule
\textbf{Name} & \textbf{train (fit)} & \textbf{test (OOS)} & \textbf{swing}\\
\midrule
HOOD & +367\% & $-$6\%  & collapse\\
MSTR & +360\% & $-$61\% & collapse\\
\midrule
MRNA & +62\%  & +103\% & \emph{improves}\\
OXY  & +27\%  & +39\%  & \emph{improves}\\
\bottomrule
\end{tabular}
\caption*{\footnotesize The optimiser found the most in-sample ``juice'' in HOOD and MSTR ---
and it was noise. On unseen data both turn negative. The genuine edges (MRNA, OXY) hold or even
strengthen out of sample.}
\end{table}

On the full-sample view HOOD looked like a strong add (194\%, Sharpe 2.57). The holdout reveals
its fitted parameters do not generalise: run forward on unseen data they lose money. Without the
out-of-sample test we would have added a name that does not work --- which is precisely why the
test is non-negotiable. MSTR is the same story, more extreme; its behaviour as a levered
Bitcoin proxy (violent gaps) breaks the limit-fill and mean-reversion assumptions the model
relies on.

% ================= fuller walk-forward =================
\section{Fuller walk-forward on the shortlist}

The 60/40 holdout is a single split, so we hardened the four diversifiers with a five-fold
expanding walk-forward: each fold \emph{re-optimises} on its train window only, freezes, and
scores the next unseen $\sim$10\% slice. Five independent out-of-sample windows per name.

\begin{table}[H]\centering\small
\renewcommand{\arraystretch}{1.2}
\begin{tabular}{l l c r r}
\toprule
\textbf{Name} & \textbf{fold OOS returns (5 unseen slices)} & \textbf{pos.\ folds} & \textbf{stitched OOS} & \textbf{ann.}\\
\midrule
\rowcolor{bckeep} DVN  & $+3 / +7 / +18 / +23 / +29$   & \textbf{5/5} & \textbf{+105\%} & 86\%\\
\rowcolor{bckeep} OXY  & $+13 / +5 / +8 / +15 / +9$    & \textbf{5/5} & +60\%  & 50\%\\
MRNA & $-10 / +13 / +35 / +33 / -1$  & 3/5 & +83\%  & 69\%\\
FSLR & $+14 / +22 / +3 / +22 / -21$  & 4/5 & +37\%  & 31\%\\
\bottomrule
\end{tabular}
\caption*{\footnotesize Return (\%) on each unseen slice, parameters re-fitted per fold on
train only. ``Stitched OOS'' compounds the five slices ($\sim$1.2y of purely out-of-sample
trading). Blue rows survive \emph{every} fold.}
\end{table}

Three things the multi-fold view adds over the single split:
\begin{itemize}[leftmargin=1.4em,itemsep=2pt]
  \item \kpi{DVN is promoted to the top.}{} The 60/40 split had shown DVN at a modest $+13\%$
        --- an artifact of which window was held out. Across five windows it is the
        \emph{strongest and most improving}: every fold positive, returns rising $+3\%\to+29\%$
        as trade counts climb ($7\to52$ buys/fold), stitched $+105\%$.
  \item \kpi{OXY is the most consistent.}{} 5/5 positive with the lowest fold-to-fold spread
        ($+5\%$ to $+15\%$) --- the dependable ballast its low-drawdown profile implied.
  \item \kpi{MRNA and FSLR are real but choppier.}{} MRNA earns its $+83\%$ from three big
        middle folds with two flat/negative ends; FSLR's most-recent fold is $-21\%$, a sign its
        edge is fading in the latest data.
\end{itemize}

% ================= recommendation =================
\section{Recommendation}

\begin{itemize}[leftmargin=1.4em,itemsep=2pt]
  \item \kpi{Add first --- survive every walk-forward fold:}{} \textbf{DVN} (energy; 5/5 folds,
        strongest and improving, superb Sharpe and low drawdown) and \textbf{OXY} (energy; 5/5
        folds, the most consistent, lowest drawdown --- the ballast name). Both decorrelated from
        the book; these are the highest-confidence adds.
  \item \kpi{Add, choppier:}{} \textbf{MRNA} (biotech; strong $+83\%$ stitched OOS but from three
        big folds with two flat ends --- real edge, lumpier ride) and \textbf{FSLR} (solar;
        decorrelated, 4/5 folds, but a $-21\%$ most-recent fold flags a fading edge --- the
        marginal name, size modestly).
  \item \kpi{Add for return, not decorrelation:}{} \textbf{SHOP} and \textbf{ARM} hold up
        strongly out of sample, but they correlate more with the book (ARM \emph{is} a
        semiconductor name). Take them if you want return, not to dilute the factor.
  \item \kpi{Size down / monitor:}{} \textbf{COIN} --- highest return and a uniquely negative
        internal hedge ($-0.07$), but a 41\% drawdown and a large in-sample-to-OOS fade; half a
        tranche at most. \textbf{RBLX} holds OOS but has effectively no internal hedge (0.85).
  \item \kpi{Reject:}{} \textbf{HOOD} and \textbf{MSTR} --- both fail the out-of-sample test.
\end{itemize}

\textbf{Best fit for the stated goal} (diluting the AI concentration): lead with \textbf{DVN and
OXY} (energy; clean 5/5 walk-forwards), then \textbf{MRNA} (biotech) and a modest \textbf{FSLR}
(solar) --- all decorrelated from the book.

\textbf{Caveats.} (1) Fitted parameters are freshly optimised; treat full-sample figures as a
ceiling and the out-of-sample columns as the trustworthy numbers. (2) All history sits in a
strong bull regime --- even a five-fold walk-forward cannot test a sustained bear market, which
these names have not seen. (3) Market cap, P/E and \$ turnover await a volume column and
fundamentals.

\vfill
{\footnotesize\color{bcgrey} Data: \texttt{Hybrid10\_Feed\_new} (OHLC only), Apr~2024--Jul~2026,
\$6M per name. Parameters fitted by robust differential evolution; out-of-sample holdout fits on
the first 60\% and scores the last 40\%. Correlations are daily close-to-close returns against
the existing ten-name book.}

\end{document}
